% sudoku.tex
% Zachary Catlin (zcatlin@purdue.edu), January 2005
%
% This file contains macros that typeset and solve (or with hard problems,
% partially solve) sudoku problems. Per the contest for which this was written,
% problem grids come in files like
% ----begin file----
% ..483..72
% .12....8.
% ..52.13..
% ....62.91
% 7..5.9..3
% 94.78....
% ..39.74..
% .5....61.
% 28..469..
% ----end file------
% i.e., nine lines of nine characters each (not counting the ending newline),
% each representing a row of the grid, with digits representing themselves and
% periods representing blank spaces.
%
% The macros make extensive use of chardef tokens as variables; they can only
% hold integers from 0 to 255, but they don't hog up count registers (the
% alternative); as the macros need to deal with large numbers of small integers,
% this works out wonderfully (thanks to Victor Eijkhout for suggesting this way
% to hold integer values in his `TeX by Topic').
%
% The actual data about the grid and its state of solution are stored in two
% sets of tokens, which I use as two- and three-dimensional arrays:
%
% \sdk@grid{x}c{y}
%    This stores the digit at row {x} and column {y} of the grid, or 0 if there
%    is no digit at that position; \sdk@input and \sdk@parseline initialize
%    these based on the given file, and \sdk@writegrid and \sdk@gridline use
%    them to typeset the grid.
%
% \sdk@q{x}c{y}n{z}
%    This stores whether it is possible, based on what is known about the grid,
%    that the square at row {x} and {y} could be the digit {z}: 1 means that
%    it is still *possible* that this could be the case, 0 means that either
%    that has been ruled out or that the digit for that square has been
%    determined. \sdk@initgrid initializes these based on \sdk@grid*c*, and
%    \sdk@solve uses them in the solution process.
%
% Both arrays are 1-based; using these as arrays means extensive use of
% \csname...\endcsname to support dynamic token names.
%
% Please note that these macros could use improvement; most notably, any
% inconsistencies in the grid will go unnoticed, and if the file is not exactly
% in the above format, Bad Things(TM) may happen. Fixing these, and adding
% improvements to the solver, are left as exercises to the reader.
%
% Thanks go to the aforementioned Victor Eijkhout and to Paul Abrahams,
% Kathryn Hargreaves, and Karl Berry for their `TeX for the Impatient'; without
% these books (and their free availability), I wouldn't have been able to do
% what I did; I heartily recommend both books to anyone who wants to learn
% how to use TeX at a low level.

% This takes chardef token #1 and increments it by #2. You can't do this with a
% primitive as chardef tokens are constants, but one can do the addition using
% a count register; \count255 is by convention used as a temporary variable,
% so it is a natural choice.
\def\sdk@addto#1#2{%
  \count255 = \number#1%
  \advance\count255 by #2%
  \chardef#1 = \count255\relax%
}

% This takes a coordinate #1 and changes it to that of the first row or column
% of the 3x3 square; specifically, this calculates ((n-1) - ((n-1) mod 3)) + 1.
\def\sdk@squarebase#1{%
  \count255 = \number#1%
  \advance\count255 by -1%
  \divide\count255 by 3%
  \multiply\count255 by -3%
  \advance\count255 by \number#1%
  \multiply\count255 by -1%
  \advance\count255 by \number#1%
  \advance\count255 by 1%
  \chardef#1 = \count255\relax%
}

% This takes the nine tokens after `\sdk@parseline' (in practice, the nine
% characters of a line in the file), assigns the \sdk@grid entries for the
% current row (0 for a . in the line; digits stand for themselves), and
% increments the row counter \sdk@row in preparation for the next line.
\def\sdk@parseline#1#2#3#4#5#6#7#8#9{%
  \ifx.#1%
    \expandafter\chardef\csname sdk@grid\number\sdk@row c1\endcsname = 0%
  \else
    \expandafter\chardef\csname sdk@grid\number\sdk@row c1\endcsname = #1%
  \fi
  \ifx.#2%
    \expandafter\chardef\csname sdk@grid\number\sdk@row c2\endcsname = 0%
  \else
    \expandafter\chardef\csname sdk@grid\number\sdk@row c2\endcsname = #2%
  \fi
  \ifx.#3%
    \expandafter\chardef\csname sdk@grid\number\sdk@row c3\endcsname = 0%
  \else
    \expandafter\chardef\csname sdk@grid\number\sdk@row c3\endcsname = #3%
  \fi
  \ifx.#4%
    \expandafter\chardef\csname sdk@grid\number\sdk@row c4\endcsname = 0%
  \else
    \expandafter\chardef\csname sdk@grid\number\sdk@row c4\endcsname = #4%
  \fi
  \ifx.#5%
    \expandafter\chardef\csname sdk@grid\number\sdk@row c5\endcsname = 0%
  \else
    \expandafter\chardef\csname sdk@grid\number\sdk@row c5\endcsname = #5%
  \fi
  \ifx.#6%
    \expandafter\chardef\csname sdk@grid\number\sdk@row c6\endcsname = 0%
  \else
    \expandafter\chardef\csname sdk@grid\number\sdk@row c6\endcsname = #6%
  \fi
  \ifx.#7%
    \expandafter\chardef\csname sdk@grid\number\sdk@row c7\endcsname = 0%
  \else
    \expandafter\chardef\csname sdk@grid\number\sdk@row c7\endcsname = #7%
  \fi
  \ifx.#8%
    \expandafter\chardef\csname sdk@grid\number\sdk@row c8\endcsname = 0%
  \else
    \expandafter\chardef\csname sdk@grid\number\sdk@row c8\endcsname = #8%
  \fi
  \ifx.#9%
    \expandafter\chardef\csname sdk@grid\number\sdk@row c9\endcsname = 0%
  \else
    \expandafter\chardef\csname sdk@grid\number\sdk@row c9\endcsname = #9%
  \fi
  \sdk@addto{\sdk@row}{1}%
}

% This allocates a file stream for reading in grid files.
\newread\sdk@file

% This opens the specified file, initializes \sdk@row for use with
% \sdk@parseline, and reads in 9 lines to be processed by \sdk@parseline.
\def\sdk@input#1{%
  \openin\sdk@file = #1%
  \chardef\sdk@row = 1%
  \loop
    \read\sdk@file to \sdk@line
    \expandafter\sdk@parseline\sdk@line
    \ifnum\sdk@row < 10%
  \repeat
  \closein\sdk@file
}

% Writes out a line of the table made by \sdk@writegrid, using the specified
% row of \sdk@grid*c*.
\def\sdk@gridline#1{&%
  \expandafter\ifnum\csname sdk@grid#1c1\endcsname < 1%
    \empty
  \else
    \expandafter\number\csname sdk@grid#1c1\endcsname
  \fi &&%
  \expandafter\ifnum\csname sdk@grid#1c2\endcsname < 1%
    \empty
  \else
    \expandafter\number\csname sdk@grid#1c2\endcsname
  \fi &&%
  \expandafter\ifnum\csname sdk@grid#1c3\endcsname < 1%
    \empty
  \else
    \expandafter\number\csname sdk@grid#1c3\endcsname
  \fi &&%
  \expandafter\ifnum\csname sdk@grid#1c4\endcsname < 1%
    \empty
  \else
    \expandafter\number\csname sdk@grid#1c4\endcsname
  \fi &&%
  \expandafter\ifnum\csname sdk@grid#1c5\endcsname < 1%
    \empty
  \else
    \expandafter\number\csname sdk@grid#1c5\endcsname
  \fi &&%
  \expandafter\ifnum\csname sdk@grid#1c6\endcsname < 1%
    \empty
  \else
    \expandafter\number\csname sdk@grid#1c6\endcsname
  \fi &&%
  \expandafter\ifnum\csname sdk@grid#1c7\endcsname < 1%
    \empty
  \else
    \expandafter\number\csname sdk@grid#1c7\endcsname
  \fi &&%
  \expandafter\ifnum\csname sdk@grid#1c8\endcsname < 1%
    \empty
  \else
    \expandafter\number\csname sdk@grid#1c8\endcsname
  \fi &&%
  \expandafter\ifnum\csname sdk@grid#1c9\endcsname < 1%
    \empty
  \else
    \expandafter\number\csname sdk@grid#1c9\endcsname
  \fi &%
  \cr
}

% Writes out a sudoku grid using the information in \sdk@grid*c*; the grid is
% exactly square, and the numbers are centered in their squares.
\def\sdk@writegrid{{%
\tt
\setbox0 = \hbox{8}%
\edef\numlen{\the\wd0}%
\dimen0 = 13.5em
\advance\dimen0 by 6.4pt % width of the vrules
\halign to \the\dimen0{\vrule width 1pt ##\tabskip=0pt plus 1fil &%
\vbox to 1.5em{\hsize=\numlen\vfil\noindent##\vfil}&\vrule ##&%
\vbox to 1.5em{\hsize=\numlen\vfil\noindent##\vfil}&\vrule ##&%
\vbox to 1.5em{\hsize=\numlen\vfil\noindent##\vfil}&\vrule width 1pt ##&%
\vbox to 1.5em{\hsize=\numlen\vfil\noindent##\vfil}&\vrule ##&%
\vbox to 1.5em{\hsize=\numlen\vfil\noindent##\vfil}&\vrule ##&%
\vbox to 1.5em{\hsize=\numlen\vfil\noindent##\vfil}&\vrule width 1pt ##&%
\vbox to 1.5em{\hsize=\numlen\vfil\noindent##\vfil}&\vrule ##&%
\vbox to 1.5em{\hsize=\numlen\vfil\noindent##\vfil}&\vrule ##&%
\vbox to 1.5em{\hsize=\numlen\vfil\noindent##\vfil}&%
\tabskip=0pt ##\vrule width 1pt \cr
\noalign{\hrule height 1pt}
\sdk@gridline{1}
\noalign{\hrule}
\sdk@gridline{2}
\noalign{\hrule}
\sdk@gridline{3}
\noalign{\hrule height 1pt}
\sdk@gridline{4}
\noalign{\hrule}
\sdk@gridline{5}
\noalign{\hrule}
\sdk@gridline{6}
\noalign{\hrule height 1pt}
\sdk@gridline{7}
\noalign{\hrule}
\sdk@gridline{8}
\noalign{\hrule}
\sdk@gridline{9}
\noalign{\hrule height 1pt}}%
}}

% Initializes \sdk@q*c*n* with 1, then sets what we know based on \sdk@grid*c*
\def\sdk@initgrid{%
  \chardef\sdk@row = 1%
  \loop
    \chardef\sdk@col = 1%
    {%
      \loop
        \chardef\sdk@num = 1%
        {%
          \loop
            \sdk@assign{\sdk@row}{\sdk@col}{\sdk@num}{1}%
            \sdk@addto{\sdk@num}{1}%
            \ifnum\sdk@num < 10%
          \repeat
        }%
        \sdk@addto{\sdk@col}{1}%
        \ifnum\sdk@col < 10%
      \repeat
    }%
    \sdk@addto{\sdk@row}{1}%
    \ifnum\sdk@row < 10%
  \repeat
  \chardef\sdk@row = 1%
  \loop
    \chardef\sdk@col = 1%
    {%
      \loop
        \expandafter\ifnum\csname sdk@grid\number\sdk@row c\number\sdk@col\endcsname > 0%
          \sdk@eliminate{\number\sdk@row}{\number\sdk@col}{\csname sdk@grid\number\sdk@row c\number\sdk@col\endcsname}%
        \fi
        \sdk@addto{\sdk@col}{1}%
        \ifnum\sdk@col < 10%
      \repeat
    }%
    \sdk@addto{\sdk@row}{1}%
    \ifnum\sdk@row < 10%
  \repeat
  \relax
}

% Sets \sdk@q{#1}c{#2}n{#3} to #4.
\def\sdk@assign#1#2#3#4{\global
  \expandafter\chardef\csname sdk@q\number #1c\number #2n\number #3\endcsname = #4}

% Gives us the token \sdk@q{#1}c{#2}n{#3}.
\def\sdk@get#1#2#3{\expandafter\number\csname sdk@q\number #1c\number #2n\number #3\endcsname}

% \sdk@eliminate and \sdk@mark are called when adding a new digit to the grid.
% \sdk@eliminate sets \sdk@q*c*n* in accordance with what we know if the number
% at row #1 and column #2 is #3: all cells in the same row and column, row and
% number, and column and number are set to 0, since none of those are possible
% then. \sdk@mark calls \sdk@eliminate, modifies \sdk@grid appropriately, and
% sets \sdk@marked to 1, indicating to \sdk@solve that a new grid square was
% solved.
\def\sdk@mark#1#2#3{%
  \global\expandafter\chardef\csname sdk@grid\number #1c\number #2\endcsname = \number #3%
  \sdk@eliminate#1#2#3%
  \global\chardef\sdk@marked = 1%
}

\def\sdk@eliminate#1#2#3{{%
    \chardef\sdk@count = 1%
    \loop
      \sdk@assign{\sdk@count}{#2}{#3}{0}%
      \sdk@assign{#1}{\sdk@count}{#3}{0}%
      \sdk@assign{#1}{#2}{\sdk@count}{0}%
      \sdk@addto{\sdk@count}{1}%
      \ifnum\sdk@count < 10%
    \repeat
    \chardef\sdk@ax = #1%
    \sdk@squarebase{\sdk@ax}%
    \chardef\sdk@ay = #2%
    \sdk@squarebase{\sdk@ay}%
    \chardef\sdk@count = \sdk@ax
    \sdk@addto{\sdk@count}{2}%
    \loop
      \chardef\sdk@county = \sdk@ay
      \sdk@addto{\sdk@county}{2}%
      {%
        \loop
          \sdk@assign{\sdk@count}{\sdk@county}{#3}{0}%
          \ifnum\sdk@county > \sdk@ay
            \sdk@addto{\sdk@county}{-1}%
        \repeat
      }%
      \ifnum\sdk@count > \sdk@ax
        \sdk@addto{\sdk@count}{-1}%
    \repeat
}}

% \sdk@elimonerow, \sdk@elimonecol, and \sdk@elimonenum implement the following
% solution technique: if there is only one square in a given row or column where
% a given number could possibly be, then that must be where it is; and if only
% one number could possibly be the occupant of a given square, than that number
% must be that occupant. The way \sdk@q*c*n* is set up, these tests are exactly
% the same, just with different indicies in different places.
\def\sdk@elimonerow{{%
  \chardef\sdk@x = 1%
  \loop
    \chardef\sdk@y = 1%
    {%
      \loop
        \chardef\sdk@z = 1%
        {%
          \chardef\sdk@sum = 0%
          \loop
            \sdk@addto{\sdk@sum}{\sdk@get{\sdk@x}{\sdk@z}{\sdk@y}}%
            \ifnum\sdk@z < 9%
              \sdk@addto{\sdk@z}{1}%
          \repeat
          \ifnum\sdk@sum = 1%
            \chardef\sdk@z = 1%
            \loop
              \expandafter\ifnum\sdk@get{\sdk@x}{\sdk@z}{\sdk@y} = 1%
                \sdk@mark{\sdk@x}{\sdk@z}{\sdk@y}%
              \fi
              \ifnum\sdk@z < 9%
                \sdk@addto{\sdk@z}{1}%
            \repeat
          \fi
        }%
        \ifnum\sdk@y < 9%
          \sdk@addto{\sdk@y}{1}%
      \repeat
    }%
    \ifnum\sdk@x < 9%
      \sdk@addto{\sdk@x}{1}%
  \repeat
}}

\def\sdk@elimonecol{{%
  \chardef\sdk@x = 1%
  \loop
    \chardef\sdk@y = 1%
    {%
      \loop
        \chardef\sdk@z = 1%
        {%
          \chardef\sdk@sum = 0%
          \loop
            \sdk@addto{\sdk@sum}{\sdk@get{\sdk@z}{\sdk@x}{\sdk@y}}%
            \ifnum\sdk@z < 9%
              \sdk@addto{\sdk@z}{1}%
          \repeat
          \ifnum\sdk@sum = 1%
            \chardef\sdk@z = 1%
            \loop
              \expandafter\ifnum\sdk@get{\sdk@z}{\sdk@x}{\sdk@y} = 1%
                \sdk@mark{\sdk@z}{\sdk@x}{\sdk@y}%
              \fi
              \ifnum\sdk@z < 9%
                \sdk@addto{\sdk@z}{1}%
            \repeat
          \fi
        }%
        \ifnum\sdk@y < 9%
          \sdk@addto{\sdk@y}{1}%
      \repeat
    }%
    \ifnum\sdk@x < 9%
      \sdk@addto{\sdk@x}{1}%
  \repeat
}}

\def\sdk@elimonenum{{%
  \chardef\sdk@x = 1%
  \loop
    \chardef\sdk@y = 1%
    {%
      \loop
        \chardef\sdk@z = 1%
        {%
          \chardef\sdk@sum = 0%
          \loop
            \sdk@addto{\sdk@sum}{\sdk@get{\sdk@x}{\sdk@y}{\sdk@z}}%
            \ifnum\sdk@z < 9%
              \sdk@addto{\sdk@z}{1}%
          \repeat
          \ifnum\sdk@sum = 1%
            \chardef\sdk@z = 1%
            \loop
              \expandafter\ifnum\sdk@get{\sdk@x}{\sdk@y}{\sdk@z} = 1%
                \sdk@mark{\sdk@x}{\sdk@y}{\sdk@z}%
              \fi
              \ifnum\sdk@z < 9%
                \sdk@addto{\sdk@z}{1}%
            \repeat
          \fi
        }%
        \ifnum\sdk@y < 9%
          \sdk@addto{\sdk@y}{1}%
      \repeat
    }%
    \ifnum\sdk@x < 9%
      \sdk@addto{\sdk@x}{1}%
  \repeat
}}

\def\sdk@elimonesquare{{%
  \chardef\sdk@x = 1%
  \loop
    \chardef\sdk@y = 1%
    {%
      \loop
        \chardef\sdk@z = 1%
        {%
          \loop
            \chardef\sdk@sum = 0%
            \chardef\sdk@bx = \sdk@x
            \chardef\sdk@by = \sdk@y
            \sdk@addto{\sdk@sum}{\sdk@get{\sdk@bx}{\sdk@by}{\sdk@z}}%
            \sdk@addto{\sdk@bx}{1}%
            \sdk@addto{\sdk@sum}{\sdk@get{\sdk@bx}{\sdk@by}{\sdk@z}}%
            \sdk@addto{\sdk@bx}{1}%
            \sdk@addto{\sdk@sum}{\sdk@get{\sdk@bx}{\sdk@by}{\sdk@z}}%
            \chardef\sdk@bx = \sdk@x
            \sdk@addto{\sdk@by}{1}%
            \sdk@addto{\sdk@sum}{\sdk@get{\sdk@bx}{\sdk@by}{\sdk@z}}%
            \sdk@addto{\sdk@bx}{1}%
            \sdk@addto{\sdk@sum}{\sdk@get{\sdk@bx}{\sdk@by}{\sdk@z}}%
            \sdk@addto{\sdk@bx}{1}%
            \sdk@addto{\sdk@sum}{\sdk@get{\sdk@bx}{\sdk@by}{\sdk@z}}%
            \chardef\sdk@bx = \sdk@x
            \sdk@addto{\sdk@by}{1}%
            \sdk@addto{\sdk@sum}{\sdk@get{\sdk@bx}{\sdk@by}{\sdk@z}}%
            \sdk@addto{\sdk@bx}{1}%
            \sdk@addto{\sdk@sum}{\sdk@get{\sdk@bx}{\sdk@by}{\sdk@z}}%
            \sdk@addto{\sdk@bx}{1}%
            \sdk@addto{\sdk@sum}{\sdk@get{\sdk@bx}{\sdk@by}{\sdk@z}}%
            \ifnum\sdk@sum = 1%
              \chardef\sdk@bx = \sdk@x
              \chardef\sdk@by = \sdk@y
              \expandafter\ifnum\sdk@get{\sdk@bx}{\sdk@by}{\sdk@z} = 1%
                \sdk@mark{\sdk@bx}{\sdk@by}{\sdk@z}%
              \fi
              \sdk@addto{\sdk@bx}{1}%
              \expandafter\ifnum\sdk@get{\sdk@bx}{\sdk@by}{\sdk@z} = 1%
                \sdk@mark{\sdk@bx}{\sdk@by}{\sdk@z}%
              \fi
              \sdk@addto{\sdk@bx}{1}%
              \expandafter\ifnum\sdk@get{\sdk@bx}{\sdk@by}{\sdk@z} = 1%
                \sdk@mark{\sdk@bx}{\sdk@by}{\sdk@z}%
              \fi
              \chardef\sdk@bx = \sdk@x
              \sdk@addto{\sdk@by}{1}%
              \expandafter\ifnum\sdk@get{\sdk@bx}{\sdk@by}{\sdk@z} = 1%
                \sdk@mark{\sdk@bx}{\sdk@by}{\sdk@z}%
              \fi
              \sdk@addto{\sdk@bx}{1}%
              \expandafter\ifnum\sdk@get{\sdk@bx}{\sdk@by}{\sdk@z} = 1%
                \sdk@mark{\sdk@bx}{\sdk@by}{\sdk@z}%
              \fi
              \sdk@addto{\sdk@bx}{1}%
              \expandafter\ifnum\sdk@get{\sdk@bx}{\sdk@by}{\sdk@z} = 1%
                \sdk@mark{\sdk@bx}{\sdk@by}{\sdk@z}%
              \fi
              \chardef\sdk@bx = \sdk@x
              \sdk@addto{\sdk@by}{1}%
              \expandafter\ifnum\sdk@get{\sdk@bx}{\sdk@by}{\sdk@z} = 1%
                \sdk@mark{\sdk@bx}{\sdk@by}{\sdk@z}%
              \fi
              \sdk@addto{\sdk@bx}{1}%
              \expandafter\ifnum\sdk@get{\sdk@bx}{\sdk@by}{\sdk@z} = 1%
                \sdk@mark{\sdk@bx}{\sdk@by}{\sdk@z}%
              \fi
              \sdk@addto{\sdk@bx}{1}%
              \expandafter\ifnum\sdk@get{\sdk@bx}{\sdk@by}{\sdk@z} = 1%
                \sdk@mark{\sdk@bx}{\sdk@by}{\sdk@z}%
              \fi
            \fi
            \ifnum\sdk@z < 9%
              \sdk@addto{\sdk@z}{1}%
          \repeat
        }%
        \ifnum\sdk@y < 7%
          \sdk@addto{\sdk@y}{3}%
      \repeat
    }%
    \ifnum\sdk@x < 7%
      \sdk@addto{\sdk@x}{3}%
  \repeat
}}

% This macro works on the following idea; keep running the various solution
% techniques (so far, only \sdk@elimonerow, \sdk@elimonecol, \sdk@elimonenum,
% though many others are possible) until no more squares' numbers get deduced
% (determined by \sdk@marked being set to one whenever \sdk@mark is called,
% i.e. whenever a solution technique determines what a square's number must be).
\def\sdk@solve{%
  \loop
    \chardef\sdk@marked = 0%
    \sdk@elimonerow\sdk@elimonecol\sdk@elimonenum
    \ifnum\sdk@marked = 1%
  \repeat
  \relax
}
